\section{Subprogram 4a: Dilatation weight \texorpdfstring{$\tfrac12$}{1/2} of the odd charges on the common UV sector}
\label{sec:subprogram4a}

In this section we determine the scaling weight of the odd charges that appear on
the Minkowski UV vacuum sector after the Subprogram~1b reconstruction and the
Subprogram~3c HLS step. The argument uses:
\begin{itemize}
  \item the indexed HLS super-Poincare seed on the UV vacuum representation
        \textnormal{(Theorem~\ref{thm:Main-HLS})},
  \item the HLS completeness statement
        \textnormal{(Theorem~\ref{thm:HLS-spinor-completeness})},
  \item and the conformal common-sector supplement
        \(\mathbf{2T}_{\mathrm{conf}}\)
        \textnormal{(Definition~\ref{def:2T-conf})}.
\end{itemize}
The logic is:
\begin{enumerate}[(1)]
  \item place the conformal even generators \(P_\mu,M_{\mu\nu},D,K_\mu\) on the
        same UV vacuum representation as the HLS odd module;
  \item use HLS completeness to prove that the dilation subgroup acts on the
        finite-dimensional HLS charge module itself;
  \item run the Jacobi argument against
        \(\{Q_\alpha^I,\bar Q_{\dot\beta J}\}\propto P_\mu\) to force
        \(\Delta(Q)=\tfrac12\).
\end{enumerate}

\subsection{Input from Subprograms 1b and 3c}

We work on the Minkowski UV chart used in Subprogram~3c. The inputs are:
\begin{itemize}
  \item the Wightman net and vacuum representation
        \((\mathcal H,\pi,U,\Omega)\) with positive-energy Poincare covariance
        from Theorem~\ref{thm:subprogram1b-Wightman};
  \item the indexed HLS odd charge module \(\mathcal V_Q\) and the operators
        \(Q_\alpha^I,\bar Q_{\dot\alpha I}\) from
        Theorem~\ref{thm:Main-HLS};
  \item the completeness statement that every translation-invariant local odd
        Weyl-spinor symmetry generator of the given chirality is already of the
        form \(Q_\alpha(v)\) for a unique \(v\in\mathcal V_Q\), from
        Theorem~\ref{thm:HLS-spinor-completeness};
  \item the same-representation conformal package
        \(\mathbf{2T}_{\mathrm{conf}}\), which places
        \(P_\mu,M_{\mu\nu},D,K_\mu\) on the same UV Hilbert space.
\end{itemize}
In this subsection we also write
\[
P_{\alpha\dot\beta}:=(\sigma^\mu)_{\alpha\dot\beta}P_\mu.
\]

\begin{lemma}[Dilatation and translations on the common UV conformal sector]
\label{thm:4a-uv-dilatation-ebc}
Assume the UV conformal common-sector interface
\(\mathbf{2T}_{\mathrm{conf}}\)
\textnormal{(Definition~\ref{def:2T-conf})} on the same vacuum representation
\((\mathcal H,\pi,\Omega)\) used in Theorem~\ref{thm:Main-HLS}. Then the
selfadjoint generators \(D\) and \(P_\mu\) act on the common Stone core supplied
by \(\mathbf{2T}_{\mathrm{conf}}\) and satisfy
\begin{equation}
\label{eq:4a-conf-relations}
  [D,P_\mu] = i P_\mu.
\end{equation}
In particular, the conformal even sector and the HLS odd charge module act on
one common UV vacuum representation with a dense common invariant core.
\end{lemma}

\begin{proof}
Definition~\ref{def:2T-conf} places the conformal generators
\(P_\mu,M_{\mu\nu},D,K_\mu\) on the common UV vacuum representation and on a
common Stone core on which \eqref{eq:4a-conf-relations} holds.
Theorem~\ref{thm:Main-HLS} places the indexed HLS odd charges on that same vacuum
representation. Intersecting the HLS invariant core with the Stone core gives the
claimed common UV vacuum representation and common invariant core.
\end{proof}

\begin{lemma}[Nontriviality of the HLS translation block]
\label{lem:4a-HLS-nonzero-translation}
Assume Theorem~\ref{thm:Main-HLS}. Then some
\(P_{\alpha\dot\beta}\) is nonzero on the common HLS core.
\end{lemma}

\begin{proof}
Let \(E:=\mathcal V_Q\). By the HLS theorem invoked in
Theorem~\ref{thm:Main-HLS}, the original nonzero odd derivation is implemented by
some nonzero \(v_0\in E\), so \(Q(v_0)\neq 0\).

Suppose for contradiction that \(P_{\alpha\dot\beta}=0\) for all
\(\alpha,\dot\beta\). Then the HLS anticommutator gives
\[
0
=
\{Q_\alpha(v_0),Q_\alpha(v_0)^\dagger\}
\]
on the common core. Hence for every core vector \(\psi\),
\[
0
=
\langle \psi,\{Q_\alpha(v_0),Q_\alpha(v_0)^\dagger\}\psi\rangle
=
\|Q_\alpha(v_0)\psi\|^2+\|Q_\alpha(v_0)^\dagger\psi\|^2.
\]
Therefore \(Q_\alpha(v_0)=0\) for both \(\alpha=1,2\), contradicting the
nontriviality of the implementing HLS charge. Thus some
\(P_{\alpha\dot\beta}\neq 0\).
\end{proof}

\begin{lemma}[Finite-dimensional odd charge module is stable under \(D\)]
\label{lem:4a-D-stability}
Assume Theorem~\ref{thm:Main-HLS},
Definition~\ref{def:2T-conf}, and
Theorem~\ref{thm:HLS-spinor-completeness}. Let \(\mathcal V_Q\) be the
finite-dimensional odd charge module furnished by
Theorem~\ref{thm:Main-HLS}, and write \(Q_\alpha(v)\) for the
corresponding odd generators. Then:
\begin{enumerate}[(a)]
  \item the dilation one-parameter subgroup acts on \(\mathcal V_Q\) through a
        finite-dimensional representation \(e^{tA_Q}\) such that on the common UV
        core
        \[
          e^{itD}Q_\alpha(v)e^{-itD}=Q_\alpha(e^{tA_Q}v);
        \]
  \item differentiating at \(t=0\) gives
        \[
          [D,Q_\alpha(v)] = i\,Q_\alpha(A_Qv);
        \]
  \item with respect to the Hermitian form determined by the indexed HLS
        anticommutator, \(A_Q+A_Q^\dagger=\mathbf 1_{\mathcal V_Q}\). In
        particular, \(A_Q-\tfrac12\mathbf 1\) is skew-Hermitian.
\end{enumerate}
\end{lemma}

\begin{proof}
Set \(E:=\mathcal V_Q\), and let
\[
\mathscr Q_+
:=
\{\text{translation-invariant local odd Weyl-spinor symmetry generators of the fixed chirality}\}.
\]
By Theorem~\ref{thm:HLS-spinor-completeness}, the map
\[
J:E\longrightarrow \mathscr Q_+,
\qquad
v\longmapsto Q(v):=(Q_1(v),Q_2(v))
\]
is a linear isomorphism. In particular, \(\mathscr Q_+\) is finite-dimensional,
and the full Weyl pair records both spinor components in a single map.

For the dilation subgroup from Definition~\ref{def:2T-conf}, define
\[
\Theta_t(X):=e^{itD}Xe^{-itD}
\qquad(X\in\mathscr Q_+).
\]
Because dilations preserve locality, translation-invariance, and Weyl-spinor
type, \(\Theta_t\) maps \(\mathscr Q_+\) to itself and satisfies
\[
\Theta_{t+s}=\Theta_t\Theta_s,
\qquad
\Theta_0=\mathbf 1_{\mathscr Q_+}.
\]
Transfer this action to \(E\) by
\[
T_t:=J^{-1}\Theta_tJ\in\operatorname{GL}(E).
\]
Then, by construction,
\[
e^{itD}Q_\alpha(v)e^{-itD}=Q_\alpha(T_tv)
\]
for both \(\alpha=1,2\).

It remains to prove continuity of \(t\mapsto T_t\). Because \(\mathscr Q_+\) is
finite-dimensional, it is enough to prove continuity of a separating family of
linear functionals. Choose a basis \(X^{(1)},\dots,X^{(m)}\) of \(\mathscr Q_+\).
By finite-dimensionality, one can choose vectors \(\phi_j,\psi_j\) in the common
UV core and local observables \(A_j\) so that the linear functionals
\[
\ell_j(X):=\langle\phi_j,X\,\pi(A_j)\Omega\rangle
\qquad(1\le j\le m)
\]
separate \(\mathscr Q_+\). For fixed \(X\in\mathscr Q_+\),
\[
\ell_j(\Theta_tX)
=
\big\langle e^{-itD}\phi_j,\,
X\,e^{-itD}\pi(A_j)\Omega\big\rangle.
\]
The conformal representation is strongly continuous, and the conformal local core
is stable under the dilation subgroup by Definition~\ref{def:2T-conf}; hence each
\(t\mapsto \ell_j(\Theta_tX)\) is continuous. Therefore \(t\mapsto\Theta_t\) is
continuous on the finite-dimensional space \(\mathscr Q_+\), and hence
\(t\mapsto T_t\) is continuous on \(E\).

The standard finite-dimensional one-parameter-group theorem therefore yields a
unique \(A_Q\in\operatorname{End}(E)\) such that
\[
T_t=e^{tA_Q}.
\]
This proves item~\textnormal{(a)}. Differentiating the defining identity at
\(t=0\) on the common UV core gives
\[
[D,Q_\alpha(v)] = i\,Q_\alpha(A_Qv),
\]
which is item~\textnormal{(b)}.

For item~\textnormal{(c)}, define the HLS Hermitian form on \(E=\mathcal V_Q\) by
\[
\{Q_\alpha(v),\bar Q_{\dot\beta}(w)\}
=
2\,\langle v,w\rangle_Q\,P_{\alpha\dot\beta}.
\]
By item~\textnormal{(b)} and selfadjointness of \(D\), the barred generators obey
\[
[D,\bar Q_{\dot\beta}(w)]
=
i\,\bar Q_{\dot\beta}(A_Q^\dagger w)
\]
under the HLS Hermitian identification. Applying \(D\) to the HLS
anticommutator and using \([D,P_{\alpha\dot\beta}]=iP_{\alpha\dot\beta}\) gives
\[
\begin{aligned}
i\,2\,\langle v,w\rangle_Q\,P_{\alpha\dot\beta}
&=
[D,\{Q_\alpha(v),\bar Q_{\dot\beta}(w)\}] \\
&=
\{[D,Q_\alpha(v)],\bar Q_{\dot\beta}(w)\}
+
\{Q_\alpha(v),[D,\bar Q_{\dot\beta}(w)]\} \\
&=
i\,2\,
\bigl(
\langle A_Qv,w\rangle_Q
+
\langle v,A_Q^\dagger w\rangle_Q
\bigr)
P_{\alpha\dot\beta}.
\end{aligned}
\]
Hence
\[
2i\,\bigl\langle (A_Q+A_Q^\dagger-\mathbf 1_E)v,w\bigr\rangle_Q\,
P_{\alpha\dot\beta}
=
0
\qquad
(v,w\in E).
\]
By Lemma~\ref{lem:4a-HLS-nonzero-translation}, choose
\((\alpha_0,\dot\beta_0)\) with \(P_{\alpha_0\dot\beta_0}\neq 0\). Specializing to
that pair, the scalar coefficient must vanish, so
\[
\bigl\langle (A_Q+A_Q^\dagger-\mathbf 1_E)v,w\bigr\rangle_Q=0
\qquad
(v,w\in E).
\]
Therefore
\[
A_Q+A_Q^\dagger=\mathbf 1_E.
\]
Equivalently, \(A_Q-\tfrac12\mathbf 1_E\) is skew-Hermitian.
\end{proof}

\begin{theorem}[Odd-charge dilatation normal form]
\label{thm:4a-Delta-half}
\label{thm:4a-D-normal-form}
Assume Lemma~\ref{lem:4a-D-stability}. Set
\[
K_Q:=A_Q-\frac12\mathbf 1_{\mathcal V_Q}.
\]
Then \(K_Q^\dagger=-K_Q\), and on the common UV core one has
\[
[D,Q_\alpha(v)]
=
\frac{i}{2}\,Q_\alpha(v)+i\,Q_\alpha(K_Qv).
\]
If the barred generators are viewed through the Hermitian identification
\(\mathcal V_Q\cong\overline{\mathcal V_Q}^{\,*}\), then likewise
\[
[D,\bar Q_{\dot\alpha}(v)]
=
\frac{i}{2}\,\bar Q_{\dot\alpha}(v)-i\,\bar Q_{\dot\alpha}(K_Qv).
\]
In particular, the only possible obstruction to exact weight~\(\tfrac12\) is a
commuting skew-Hermitian endomorphism of the internal multiplicity space.
\end{theorem}

\begin{proof}
Lemma~\ref{lem:4a-D-stability} gives
\[
A_Q+A_Q^\dagger=\mathbf 1_{\mathcal V_Q},
\]
hence \(K_Q^\dagger=-K_Q\). Substituting
\(A_Q=\frac12\mathbf 1_{\mathcal V_Q}+K_Q\) into the commutator formula of
Lemma~\ref{lem:4a-D-stability} immediately gives
\[
[D,Q_\alpha(v)]
=
i\,Q_\alpha(A_Qv)
=
\frac{i}{2}\,Q_\alpha(v)+i\,Q_\alpha(K_Qv).
\]
For the barred sector, take Hilbert adjoints and use the selfadjointness of \(D\):
\[
[D,\bar Q_{\dot\alpha}(v)]
=
-[D,Q_\alpha(v)]^\dagger
=
\frac{i}{2}\,\bar Q_{\dot\alpha}(v)-i\,\bar Q_{\dot\alpha}(K_Qv).
\]
No further conclusion about \(K_Q\) follows from the stated interfaces alone.
\end{proof}

\begin{remark}[Meaning of the \(K_Q\) term]
Theorem~\ref{thm:4a-Delta-half} isolates the most general dilatation normal form
allowed by the stated interfaces. The endomorphism \(K_Q\) measures the
internal skew-Hermitian part of the dilatation action on the multiplicity
space. In Subprogram~4b that term is absorbed into the internal degree-zero
sector once that sector is realized on the same Hilbert space.
\end{remark}

\begin{corollary}[Standard half-weight after central internal normalization]
\label{cor:4b-dilatation-normalization}
Assume Theorem~\ref{thm:4a-Delta-half},
Theorem~\ref{thm:QS-conformal-plus-internal},
Theorem~\ref{thm:internal-action-from-mixed-brackets},
Lemma~\ref{lem:4b-degree-zero-faithfulness},
Lemma~\ref{lem:4b-KQ-scalar},
Lemma~\ref{lem:4b-upper-D-normal-form}, and
Corollary~\ref{cor:4b-contragredient-internal-action}. Let
\[
K_Q=\kappa\,\mathbf 1_E,
\qquad
R_{K_Q}:=T_{K_Q},
\qquad
D_{\mathrm{std}}:=D-iR_{K_Q}.
\]
Then
\[
[D_{\mathrm{std}},Q_\alpha(v)]
=
\frac{i}{2}\,Q_\alpha(v),
\qquad
[D_{\mathrm{std}},\bar Q_{\dot\alpha}(v)]
=
\frac{i}{2}\,\bar Q_{\dot\alpha}(v).
\]
\[
[D_{\mathrm{std}},S^\alpha(\lambda)]
=
-\frac{i}{2}\,S^\alpha(\lambda),
\qquad
[D_{\mathrm{std}},\bar S^{\dot\alpha}(\lambda)]
=
-\frac{i}{2}\,\bar S^{\dot\alpha}(\lambda).
\]
Moreover, for every \(A\in\operatorname{End}(E)\),
\[
  [D_{\mathrm{std}},T_A]=0.
\]
Finally,
\[
  [D_{\mathrm{std}},P_\mu]=iP_\mu,
  \qquad
  [D_{\mathrm{std}},M_{\mu\nu}]=[D,M_{\mu\nu}],
\]
\end{corollary}

\begin{proof}
Choose an orthonormal basis \(\{e_I\}_{I=1}^N\) of \(E\) with dual basis
\(\{e^I\}_{I=1}^N\), and write
\[
R^I{}_J:=T_{e_I,e^J},
\qquad
r:=\sum_I R^I{}_I.
\]
Since \(K_Q=\kappa\,\mathbf 1_E\), one has
\[
R_{K_Q}=\kappa r.
\]
By Corollary~\ref{cor:4b-contragredient-internal-action},
\[
\bigl[r,Q_\alpha^I\bigr]=Q_\alpha^I,
\qquad
\bigl[r,\bar Q_{\dot\alpha I}\bigr]=-\bar Q_{\dot\alpha I},
\]
\[
\bigl[r,S^\alpha{}_I\bigr]=-S^\alpha{}_I,
\qquad
\bigl[r,\bar S^{\dot\alpha I}\bigr]=\bar S^{\dot\alpha I}.
\]
Hence
\[
\bigl[iR_{K_Q},Q_\alpha^I\bigr]=i\kappa\,Q_\alpha^I,
\qquad
\bigl[iR_{K_Q},\bar Q_{\dot\alpha I}\bigr]=-i\kappa\,\bar Q_{\dot\alpha I},
\]
\[
\bigl[iR_{K_Q},S^\alpha{}_I\bigr]=-i\kappa\,S^\alpha{}_I,
\qquad
\bigl[iR_{K_Q},\bar S^{\dot\alpha I}\bigr]=i\kappa\,\bar S^{\dot\alpha I}.
\]

On the other hand, Theorem~\ref{thm:4a-Delta-half},
Lemma~\ref{lem:4b-KQ-scalar}, and
Lemma~\ref{lem:4b-upper-D-normal-form} give
\[
[D,Q_\alpha^I]
=
\Bigl(\frac{i}{2}+i\kappa\Bigr)Q_\alpha^I,
\qquad
[D,\bar Q_{\dot\alpha I}]
=
\Bigl(\frac{i}{2}-i\kappa\Bigr)\bar Q_{\dot\alpha I},
\]
\[
[D,S^\alpha{}_I]
=
\Bigl(-\frac{i}{2}-i\kappa\Bigr)S^\alpha{}_I,
\qquad
[D,\bar S^{\dot\alpha I}]
=
\Bigl(-\frac{i}{2}+i\kappa\Bigr)\bar S^{\dot\alpha I}.
\]
Subtracting the \(iR_{K_Q}\)-commutators therefore yields the exact half-weight
relations for all four odd families.

For the internal block, let \(A\in\operatorname{End}(E)\) and set
\[
  U_A:=[R_{K_Q},T_A].
\]
Because \(R_{K_Q}\) and \(T_A\) are degree-zero internal operators,
\(U_A\) is again degree-zero internal. By linearity from
Theorem~\ref{thm:QS-conformal-plus-internal},
\[
[D,T_A]=0,
\qquad
[T_A,P_\mu]=[T_A,M_{\mu\nu}]=0.
\]
Moreover, for every \(w\in E\),
\[
\begin{aligned}
[U_A,Q_\alpha(w)]
&=
[R_{K_Q},Q_\alpha(Aw)]-[T_A,Q_\alpha(K_Qw)]\\
&=
Q_\alpha(K_QAw)-Q_\alpha(AK_Qw)\\
&= 
0
\end{aligned}
\]
by Lemma~\ref{lem:4b-KQ-scalar}. Hence \(U_A\) is a degree-zero internal
operator acting trivially on the lower odd layer. By
Lemma~\ref{lem:4b-degree-zero-faithfulness}, it acts trivially on the entire odd
closure, so \(U_A=0\). Therefore
\[
[R_{K_Q},T_A]=0
\qquad
(A\in\operatorname{End}(E)),
\]
and thus
\[
  [D_{\mathrm{std}},T_A]
  =
  [D,T_A]-i[R_{K_Q},T_A]
  =
  0.
\]
Since \(R_{K_Q}\) is internal, it already commutes with \(P_\mu\) and
\(M_{\mu\nu}\), so
\[
[D_{\mathrm{std}},P_\mu]=[D,P_\mu]=iP_\mu,
\qquad
[D_{\mathrm{std}},M_{\mu\nu}]=[D,M_{\mu\nu}].
\]
\end{proof}

\begin{remark}[Current-theoretic consequence of \(D_{\mathrm{std}}\)]
For any conserved local current representative of an HLS charge that is
conformal-primary with respect to \(D_{\mathrm{std}}\), the relations of
Corollary~\ref{cor:4b-dilatation-normalization} force engineering dimension
\(\tfrac72\).
\end{remark}


\section{Subprogram 4b: Same-\texorpdfstring{$\mathcal H$}{H} conformal closure of the HLS odd sector}
\label{sec:subprogram4b}

In this section the HLS odd sector is closed conformally on the same UV Hilbert
space. The argument proceeds as follows:
\begin{enumerate}[(1)]
  \item strengthen the HLS import once, by using
        Theorem~\ref{thm:HLS-spinor-completeness};
  \item take the finite conformal closure of the HLS odd seed;
  \item prove same-\(\mathcal H\) implementability of that finite closure on a
        common invariant core;
  \item recover the internal degree-zero algebra from mixed \(Q\)-\(S\) brackets
        and prove saturation to \(\mathfrak u(N)\);
  \item identify the resulting operator superalgebra with the standard \(4\)D
        superconformal algebra.
\end{enumerate}
Theorem~\ref{thm:indexed-current-reconstruction} from the main spine gives a
separate AQFT result for choosing indexed current representatives, while the
closure argument below is formulated directly at the
operator and module level.

\subsection{Representation-theoretic input and odd conformal closure}
\label{subsec:4b-closure}

\begin{theorem}[Generated conformal prolongation of the HLS seed
  \cite{Kac1977,Nahm:1977tg}]
\label{thm:4b-odd-prolongation}
Assume Theorem~\ref{thm:Main-HLS},
Definition~\ref{def:2T-conf},
Theorem~\ref{thm:HLS-spinor-completeness}, and
Theorem~\ref{thm:4a-Delta-half}. Let \(E:=\mathcal V_Q\), and let \(W_Q\) be the
translation-invariant HLS odd seed, viewed through the normal form of
Theorem~\ref{thm:4a-Delta-half} as the tensor product of the standard
weight-\(\tfrac12\) Weyl module with the internal multiplicity space \(E\). Let
\[
W_{\mathrm{sc}}
:=
U\!\bigl(\mathfrak{so}(4,2)_\mathbb C\bigr)\cdot W_Q
\]
inside the odd local-derivation space on the UV vacuum representation. Then:
\begin{enumerate}[(i)]
  \item \(W_{\mathrm{sc}}\) is finite-dimensional.
  \item \(W_{\mathrm{sc}}\) has exactly two odd weight layers,
\[
W_{\mathrm{sc}} = W_Q \oplus W_S.
\]
  \item with respect to that normalized conformal part, \(W_Q\) has weight
        \(+\tfrac12\) and \(W_S\) has weight \(-\tfrac12\).
  \item \(P_\mu\) annihilates \(W_Q\), and \(K_\mu\) annihilates \(W_S\).
  \item No further odd layers occur.
\end{enumerate}
\end{theorem}

\begin{proof}
Let
\[
  \mathfrak g:=\mathfrak{so}(4,2)_\mathbb C
  =
  \mathfrak g_{-1}\oplus\mathfrak g_0\oplus\mathfrak g_{+1}
\]
be the standard conformal \(|1|\)-grading, with
\[
  \mathfrak g_{+1}\cong\operatorname{span}\{P_\mu\},
  \qquad
  \mathfrak g_{-1}\cong\operatorname{span}\{K_\mu\},
  \qquad
  \mathfrak g_0\cong\mathfrak{so}(1,3)_\mathbb C\oplus\mathbb C D.
\]
By Theorem~\ref{thm:Main-HLS} and Theorem~\ref{thm:4a-Delta-half}, the HLS seed
\(W_Q\) carries the Weyl-spinor factor with standard weight \(\tfrac12\), while the
skew-Hermitian part \(K_Q\) acts only on the multiplicity space \(E\) and commutes
with the Lorentz-spinor factor. Hence \(W_Q\) may be viewed as a
finite-dimensional \(\mathfrak g_0\)-module whose conformal part is the standard
lowest-weight Weyl module, annihilated by \(\mathfrak g_{+1}\) because the HLS
generators are translation-invariant. Thus \(W_Q\) is naturally a module for the
parabolic subalgebra
\[
  \mathfrak p^+:=\mathfrak g_0\oplus\mathfrak g_{+1}.
\]
Form the generalized Verma module
\[
  M(W_Q):=U(\mathfrak g)\otimes_{U(\mathfrak p^+)}W_Q.
\]
The concrete conformal closure in Definition~\ref{def:odd-conformal-closure} is
exactly the image of \(M(W_Q)\) in the odd local-derivation space on the UV
vacuum representation.

Chirality by chirality, the conformal part of \(W_Q\) is \(E\) tensored with the
lowest-weight Weyl spinor of weight \(\tfrac12\). The imported conformal-spinor
prolongation classification cited in the appendix (\cite{Kac1977,Nahm:1977tg})
says that the generated module for such a seed has the standard two-layer
conformal-spinor quotient and no higher odd prolongation. Therefore the concrete
image generated by the actual HLS seed factors through a finite-dimensional module
with exactly the two odd normalized weights \(\pm\tfrac12\). The lowest layer is
the seed module \(W_Q\), the upper layer is its conformal prolongation \(W_S\),
\(P_\mu\) kills the lower layer, \(K_\mu\) kills the upper layer, and no further
odd layers occur.
\end{proof}

\begin{definition}[Odd conformal closure]
\label{def:odd-conformal-closure}
Let \(W_Q\) be the odd derivation module generated by the HLS charges of
Theorem~\ref{thm:Main-HLS}; equivalently,
\[
W_Q
:=
\left\{
A\mapsto i[Q(v),\pi(A)]_{\mathrm s}
\;:\;
v\in\mathcal V_Q
\right\}.
\]
Define \(W_{\mathrm{sc}}\) to be the smallest linear space of odd local
derivations on the UV net that:
\begin{itemize}
  \item contains \(W_Q\), and
  \item is stable under commutator with the even conformal generators
        \(P_\mu,M_{\mu\nu},D,K_\mu\).
\end{itemize}
We call \(W_{\mathrm{sc}}\) the post-HLS odd conformal closure of the HLS seed.
\end{definition}

\begin{theorem}[Finite-dimensional conformal closure of the HLS odd sector]
\label{thm:odd-conformal-closure-finite}
Assume Theorem~\ref{thm:4b-odd-prolongation}. Then \(W_{\mathrm{sc}}\) is
finite-dimensional and splits as
\[
W_{\mathrm{sc}} = W_Q \oplus W_S,
\]
where:
\begin{itemize}
  \item with respect to the normalized conformal part, \(W_Q\) has weight
        \(+\tfrac12\),
  \item with respect to the normalized conformal part, \(W_S\) has weight
        \(-\tfrac12\),
  \item \(P_\mu\) annihilates \(W_Q\),
  \item \(K_\mu\) annihilates \(W_S\),
  \item and \(W_S\) has exactly the left/right Weyl-spinor chirality pattern of
        the standard \((S,\bar S)\) layer.
\end{itemize}
\end{theorem}

\begin{proof}
This is the content of Theorem~\ref{thm:4b-odd-prolongation}, restated with the
notation of Definition~\ref{def:odd-conformal-closure}.
\end{proof}

\subsection{Same-\texorpdfstring{$\mathcal H$}{H} implementability}
\label{subsec:4b-sameH}

\begin{lemma}[Conformal Nelson smooth domain]
\label{lem:4b-conformal-nelson-domain}
Let
\[
N_{\mathrm{conf}}
:=
1+\sum_\mu P_\mu^2+\sum_{\mu<\nu}M_{\mu\nu}^2+D^2+\sum_\mu K_\mu^2
\]
with closure taken on the common Stone core from
Definition~\ref{def:2T-conf}. Then
\[
\mathcal D_\infty(N_{\mathrm{conf}})
:=
\bigcap_{r\geq 0}\operatorname{Dom}N_{\mathrm{conf}}^r
\]
is dense, invariant under the connected conformal group, and coincides with the
smooth-vector space for the conformal representation.
\end{lemma}

\begin{proof}
This is the standard Nelson theorem for a finite family of essentially
selfadjoint generators of a strongly continuous unitary representation of a
finite-dimensional Lie group. Definition~\ref{def:2T-conf} supplies exactly those
hypotheses for the conformal generator family
\(\{P_\mu,M_{\mu\nu},D,K_\mu\}\).
\end{proof}

\begin{lemma}[The indexed HLS seed has a common linear \(H\)-bound]
\label{lem:4b-HLS-seed-H-bound}
Assume Theorem~\ref{thm:Main-HLS}. Then, for any chosen norm on the two-component
Weyl pair \(Q(v)=(Q_1(v),Q_2(v))\), there exists \(C_0>0\) such that
\[
\|Q(v)\psi\|
\le
C_0\,\|v\|_Q\,\|(1+H)\psi\|
\qquad
(v\in\mathcal V_Q,\ \psi\in\mathcal D_Q),
\]
where \(H:=P^0\). In particular, the seed polynomial energy bound holds
with \(m_0=1\).
\end{lemma}

\begin{proof}
Choose an orthonormal basis \(\{e_I\}_{I=1}^N\) of \(E:=\mathcal V_Q\), write
\[
v=\sum_{I=1}^N v_I e_I,
\]
and define
\[
Q_\alpha(v):=\sum_{I=1}^N v_I Q_\alpha^I,
\qquad
\bar Q_{\dot\beta}(v):=\sum_{I=1}^N \overline{v_I}\,\bar Q_{\dot\beta I}.
\]
By bilinearity and the indexed HLS anticommutator from
Theorem~\ref{thm:Main-HLS},
\[
\{Q_\alpha(v),\bar Q_{\dot\beta}(v)\}
=
2\,\|v\|_Q^2\,\sigma^\mu_{\alpha\dot\beta}P_\mu
\]
on the common invariant core \(\mathcal D_Q\). Summing the diagonal entries
\(\alpha=\dot\alpha=1,2\) and using \(\sigma^0=\mathbf 1_2\) together with the
tracelessness of the spatial Pauli matrices gives
\[
\sum_{\alpha=1}^2 \{Q_\alpha(v),\bar Q_{\dot\alpha}(v)\}
=
4\,\|v\|_Q^2\,H.
\]
Therefore, for every \(\psi\in\mathcal D_Q\),
\[
\sum_{\alpha=1}^2
\Bigl(
\|Q_\alpha(v)\psi\|^2+\|\bar Q_{\dot\alpha}(v)\psi\|^2
\Bigr)
=
4\,\|v\|_Q^2\,\langle\psi,H\psi\rangle.
\]
Because Theorem~\ref{thm:Main-HLS} is built on the positive-energy vacuum
representation from Theorem~\ref{thm:subprogram1b-Wightman}, the spectrum condition
gives \(H\geq 0\). Hence the spectral theorem yields
\[
\langle\psi,H\psi\rangle
\le
\|(1+H)\psi\|^2.
\]
It follows that
\[
\sum_{\alpha=1}^2 \|Q_\alpha(v)\psi\|^2
\le
4\,\|v\|_Q^2\,\|(1+H)\psi\|^2.
\]
If \(\|\cdot\|_{\mathrm W}\) denotes the chosen norm on the two-component Weyl pair,
finite-dimensional norm equivalence gives \(c_{\mathrm W}>0\) such that for every
\((\xi_1,\xi_2)\in\mathcal H\oplus\mathcal H\),
\[
\|(\xi_1,\xi_2)\|_{\mathrm W}
\le
c_{\mathrm W}\,(\|\xi_1\|^2+\|\xi_2\|^2)^{1/2}
\]
Applying this pointwise to
\((Q_1(v)\psi,Q_2(v)\psi)\) yields
\[
\|Q(v)\psi\|
\le
2c_{\mathrm W}\,\|v\|_Q\,\|(1+H)\psi\|.
\]
Thus the stated bound holds with \(C_0:=2c_{\mathrm W}\), and in particular one may
take \(m_0=1\).
\end{proof}

\begin{lemma}[Seed bounds transfer from \(H\) to \(N_{\mathrm{conf}}\)]
\label{lem:4b-seed-bounds-to-nconf}
Assume Theorem~\ref{thm:Main-HLS} and Definition~\ref{def:2T-conf}. Then, after
restricting to
\(\mathcal D_\infty(N_{\mathrm{conf}})\cap\mathcal D_Q\), there exist \(m_1\geq 0\)
and \(C_1>0\) such that
\[
\|Q(v)\psi\|
\leq
C_1\,\|v\|_Q\,\|(1+N_{\mathrm{conf}})^{m_1}\psi\|
\qquad
(v\in\mathcal V_Q,\ \psi\in\mathcal D_\infty(N_{\mathrm{conf}})\cap\mathcal D_Q).
\]
\end{lemma}

\begin{proof}
By Lemma~\ref{lem:4b-HLS-seed-H-bound}, there exists \(C_0>0\) such that
\[
\|Q(v)\psi\|
\le
C_0\,\|v\|_Q\,\|(1+H)\psi\|
\qquad
(v\in\mathcal V_Q,\ \psi\in\mathcal D_Q).
\]
Since \(H=P^0\) is one of the conformal generators entering \(N_{\mathrm{conf}}\),
the conformal Nelson domain carries a polynomial comparison
\[
  \|(1+H)\psi\|
  \leq
  C\,\|(1+N_{\mathrm{conf}})^{m_1}\psi\|
\]
for suitable \(m_1\) and \(C\). Combining the two inequalities gives the claim
with \(C_1:=CC_0\).
\end{proof}

\begin{lemma}[One conformal derivative costs one Sobolev order]
\label{lem:4b-one-conformal-derivative}
Let \(X\) be a closable operator satisfying
\[
\|X\psi\|
\leq
C\,\|(1+N_{\mathrm{conf}})^m\psi\|
\qquad
(\psi\in\mathcal D_\infty(N_{\mathrm{conf}})).
\]
For any one-parameter conformal subgroup with selfadjoint generator \(G\), the
orbit map
\[
t\longmapsto e^{itG}Xe^{-itG}
\]
is differentiable as a map
\[
\operatorname{Dom}(1+N_{\mathrm{conf}})^{(m+1)/2}\longrightarrow\mathcal H,
\]
and its derivative at \(t=0\) is the commutator \([iG,X]\) on
\(\mathcal D_\infty(N_{\mathrm{conf}})\).
\end{lemma}

\begin{proof}
The conformal group acts continuously on each Nelson-Sobolev space built from
\(N_{\mathrm{conf}}\). Differentiating the conjugation orbit loses one group
derivative and hence one Sobolev order. This is the standard smooth-vector
calculus for finite-dimensional unitary Lie-group representations.
\end{proof}

\begin{theorem}[Uniform orbit energy bounds for the odd conformal closure]
\label{thm:conformal-orbit-uniform-energy-bounds}
Assume Theorem~\ref{thm:Main-HLS},
Definition~\ref{def:2T-conf}, and
Theorem~\ref{thm:odd-conformal-closure-finite}.
Then there exist an integer \(M\geq 0\), a constant \(C>0\), and a norm
\(\|\cdot\|_{\mathrm{sc}}\) on \(W_{\mathrm{sc}}\) such that every odd derivation
in \(W_{\mathrm{sc}}\) is implemented on
\(\mathcal D_\infty(N_{\mathrm{conf}})\cap\mathcal D_Q\) by a closable operator
\(X\) satisfying
\[
\|X\psi\|
\le
C\,\|X\|_{\mathrm{sc}}\,\|(1+N_{\mathrm{conf}})^M\psi\|
\qquad
(\psi\in\mathcal D_\infty(N_{\mathrm{conf}})\cap\mathcal D_Q).
\]
In particular, the entire odd conformal closure obeys one common polynomial graph
bound on the same UV Hilbert space.
\end{theorem}

\begin{proof}
By Lemma~\ref{lem:4b-seed-bounds-to-nconf}, the HLS seed implementers satisfy a
common \(N_{\mathrm{conf}}\)-Sobolev bound on
\(\mathcal D_\infty(N_{\mathrm{conf}})\cap\mathcal D_Q\). By
Theorem~\ref{thm:odd-conformal-closure-finite}, the whole odd conformal closure
has depth one:
\[
  W_{\mathrm{sc}}=W_Q\oplus W_S,
  \qquad
  W_S=[K,W_Q].
\]
The \(W_Q\)-layer already satisfies the transferred Nelson bound. Every element of
\(W_S\) is obtained by one conformal derivative of a seed implementer, so
Lemma~\ref{lem:4b-one-conformal-derivative} yields the same kind of bound with one
additional Sobolev order. Because \(W_{\mathrm{sc}}\) is finite-dimensional, one
may choose a basis \(\{X_a\}\) of \(W_Q\oplus W_S\), let \(M\) be the maximum of
the finitely many required Sobolev exponents, and define
\[
  \|X\|_{\mathrm{sc}}
  :=
  \sum_a |c_a|
  \qquad
  \text{for }
  X=\sum_a c_a X_a.
\]
Then one constant \(C\) works for every \(X\in W_{\mathrm{sc}}\), giving the
displayed bound.
\end{proof}

\begin{lemma}[G\r{a}rding vectors from a dense Hilbert subspace are smooth-dense]
\label{lem:4b-garding-dense-from-dense-set}
Let \(G\) be a finite-dimensional Lie group with strongly continuous unitary
representation \(U\) on \(\mathcal H\), and let \(\mathcal H^\infty\) be the
smooth-vector Fr\'echet space. If \(E\subset\mathcal H\) is dense, then
\[
\mathcal G(E)
:=
\mathrm{span}\{U(f)\eta:f\in C_c^\infty(G),\ \eta\in E\}
\]
is Fr\'echet-dense in \(\mathcal H^\infty\).
\end{lemma}

\begin{proof}
Fix \(\psi\in\mathcal H^\infty\) and a finite family of Fr\'echet seminorms
\(p_{Y_1,\dots,Y_m}(\xi):=\|dU(Y_1)\cdots dU(Y_m)\xi\|\). Let \((f_n)\) be a
standard approximate identity in \(C_c^\infty(G)\). Then
\[
U(f_n)\psi\longrightarrow \psi
\]
in every smooth-vector seminorm. For each fixed \(n\), the operators
\(U(L_{Y_1}\cdots L_{Y_m}f_n)\) are bounded on \(\mathcal H\). Since \(E\) is dense in
\(\mathcal H\), one may choose \(\eta_n\in E\) such that
\[
\|U(L_{Y_1}\cdots L_{Y_m}f_n)(\eta_n-\psi)\|
\]
is arbitrarily small for the finitely many seminorms under consideration. Then
\[
p_{Y_1,\dots,Y_m}\bigl(U(f_n)\eta_n-\psi\bigr)
\le
p_{Y_1,\dots,Y_m}\bigl(U(f_n)(\eta_n-\psi)\bigr)
+
p_{Y_1,\dots,Y_m}\bigl(U(f_n)\psi-\psi\bigr)
\]
can be made arbitrarily small. A diagonal argument over a countable defining
family of seminorms shows that \(\mathcal G(E)\) is Fr\'echet-dense in
\(\mathcal H^\infty\).
\end{proof}

\begin{lemma}[Finite conformal operator families preserve conformal smooth vectors]
\label{lem:4b-smooth-vector-transfer}
Assume Definition~\ref{def:2T-conf},
Theorem~\ref{thm:conformal-orbit-uniform-energy-bounds}, and that
\(W_{\mathrm{sc}}\) is finite-dimensional and stable under the even conformal
action. Set
\[
\mathcal D_{\mathrm{seed}}
:=
\mathcal D_\infty(N_{\mathrm{conf}})\cap\mathcal D_Q.
\]
Then every \(X\in W_{\mathrm{sc}}\) extends uniquely from
\(\mathcal D_{\mathrm{seed}}\) to a continuous endomorphism of the Fr\'echet
space \(\mathcal D_\infty(N_{\mathrm{conf}})\), and the whole family
\(W_{\mathrm{sc}}\) preserves \(\mathcal D_\infty(N_{\mathrm{conf}})\).
\end{lemma}

\begin{proof}
By Lemma~\ref{lem:4b-conformal-nelson-domain},
\(\mathcal D_\infty(N_{\mathrm{conf}})\) is dense in \(\mathcal H\), and by
Theorem~\ref{thm:Main-HLS} the HLS core \(\mathcal D_Q\) is dense in \(\mathcal H\).
Hence their intersection \(\mathcal D_{\mathrm{seed}}\) is dense in \(\mathcal H\).

Let \(G:=\mathrm{Conf}_0(\mathbb R^{1,3})\), let \(U:=U_{\mathrm{conf}}\), and for
\(X\in W_{\mathrm{sc}}\) write
\[
\Theta_g(X):=U(g)XU(g)^{-1}.
\]
Choose a basis \(X_1,\dots,X_r\) of the finite-dimensional space \(W_{\mathrm{sc}}\)
and write
\[
\Theta_g(X)=\sum_{a=1}^r c_a(g;X)X_a
\]
with smooth coefficient functions \(c_a(g;X)\).

For \(f\in C_c^\infty(G)\) and \(\psi\in\mathcal D_{\mathrm{seed}}\), define
\[
X\,U(f)\psi
:=
\int_G f(g)\,U(g)\,\Theta_{g^{-1}}(X)\psi\,dg.
\]
Because \(\Theta_{g^{-1}}(X)\) stays in the finite-dimensional family
\(W_{\mathrm{sc}}\), the uniform Sobolev bound from
Theorem~\ref{thm:conformal-orbit-uniform-energy-bounds} implies that the
integrand is norm-continuous and compactly supported in \(g\); thus the integral
is well defined. Expanding \(\Theta_{g^{-1}}(X)\) in the basis gives
\[
XU(f)\psi
=
\sum_{a=1}^r\int_G f(g)c_a(g^{-1};X)U(g)X_a\psi\,dg,
\]
so \(XU(f)\psi\) is a finite sum of conformal smooth vectors.

Now apply any Lie derivative \(dU(Y_1)\cdots dU(Y_m)\) to the above formula.
Differentiation under the integral is allowed because orbit derivatives of
\(\Theta_g(X)\) still lie in \(W_{\mathrm{sc}}\), and each such derivative obeys
the same kind of polynomial Sobolev bound by repeated use of
Lemma~\ref{lem:4b-one-conformal-derivative} and finite-dimensionality. Hence, for
each Fr\'echet seminorm \(p_m\) on \(\mathcal D_\infty(N_{\mathrm{conf}})\), there
exist seminorms \(q_{m,X}\) on \(\mathcal D_{\mathrm{seed}}\) and constants
\(C_{m,X,f}\) such that
\[
p_m\bigl(XU(f)\psi\bigr)
\le
C_{m,X,f}\,q_{m,X}(\psi).
\]
Therefore \(X\) acts continuously on the G\r{a}rding subspace generated from
\(\mathcal D_{\mathrm{seed}}\).

By Lemma~\ref{lem:4b-garding-dense-from-dense-set}, the conformal G\r{a}rding
vectors generated from the dense set \(\mathcal D_{\mathrm{seed}}\) are
Fr\'echet-dense in
\(\mathcal D_\infty(N_{\mathrm{conf}})\). The seminorm estimates above therefore
extend \(X\) uniquely and continuously to all conformal smooth vectors. This
extension preserves \(\mathcal D_\infty(N_{\mathrm{conf}})\).
\end{proof}

\begin{theorem}[Same-UV-Hilbert-space odd conformal closure package]
\label{thm:4b-narrow-common-sector}
\label{thm:odd-conformal-closure-implementable}
Assume Theorem~\ref{thm:Main-HLS},
Definition~\ref{def:2T-conf},
Theorem~\ref{thm:4a-Delta-half},
Theorem~\ref{thm:odd-conformal-closure-finite}, and
Theorem~\ref{thm:conformal-orbit-uniform-energy-bounds}. Then every derivation in
\(W_{\mathrm{sc}}\) is implemented by a closable operator on the same UV Hilbert
space \(\mathcal H\), and the dense domain
\[
\mathcal D_{\mathrm{sc}}
:=
\mathcal D_\infty(N_{\mathrm{conf}})
\subset\mathcal H
\]
is a common invariant core for
\[
P_\mu,\ M_{\mu\nu},\ D,\ K_\mu
\]
and for all odd implementers belonging to \(W_{\mathrm{sc}}\). After choosing
bases of the two odd layers furnished by Theorem~\ref{thm:4b-odd-prolongation},
one obtains operators
\[
Q_\alpha^I,\ \bar Q_{\dot\alpha I},\ S^\alpha{}_I,\ \bar S^{\dot\alpha I}
\]
on that same common core.
\end{theorem}

\begin{proof}
Lemma~\ref{lem:4b-conformal-nelson-domain} gives a dense conformal smooth-vector
space \(\mathcal D_{\mathrm{sc}}\) preserved by the even conformal generators.
Theorem~\ref{thm:conformal-orbit-uniform-energy-bounds} gives one common
polynomial \(N_{\mathrm{conf}}\)-Sobolev bound for all
\(X\in W_{\mathrm{sc}}\) on the seed domain
\[
\mathcal D_{\mathrm{seed}}
:=
\mathcal D_\infty(N_{\mathrm{conf}})\cap\mathcal D_Q.
\]
Lemma~\ref{lem:4b-smooth-vector-transfer} upgrades each
\(X\in W_{\mathrm{sc}}\) to a continuous endomorphism of
\(\mathcal D_\infty(N_{\mathrm{conf}})\). In particular, every odd derivation in
\(W_{\mathrm{sc}}\) is implemented by a closable operator on the same UV Hilbert
space and preserves \(\mathcal D_\infty(N_{\mathrm{conf}})\).

Therefore \(\mathcal D_{\mathrm{sc}}=\mathcal D_\infty(N_{\mathrm{conf}})\) is a
dense common invariant core for the full even conformal sector and the entire odd
conformal closure. Naming the lower and upper odd layers by
\(Q_\alpha^I,\bar Q_{\dot\alpha I}\) and \(S^\alpha{}_I,\bar S^{\dot\alpha I}\)
gives the stated same-\(\mathcal H\) realization.
\end{proof}

\subsection{Mixed brackets and the internal algebra}
\label{subsec:4b-internal}

Write \(E:=\mathcal V_Q\), with \(\dim_\mathbb C E=N\). The HLS Hermitian form
from Theorem~\ref{thm:Main-HLS} identifies \(E\) anti-linearly with \(E^*\); write
\[
  \flat:E\to E^*,
  \quad
  v\mapsto v^\flat,
  \qquad
  \sharp:E^*\to E,
  \quad
  \lambda\mapsto \lambda^\sharp,
\]
so that \(\lambda(v)=\langle v,\lambda^\sharp\rangle_Q\). For
\(A\in\operatorname{End}(E)\), write
\[
  A^\vee:E^*\to E^*,
  \qquad
  (A^\vee\lambda)(v):=\lambda(Av),
\]
and let \(A^\star\in\operatorname{End}(E)\) be defined by
\[
  \langle Av,w\rangle_Q=\langle v,A^\star w\rangle_Q.
\]
For rank-one operators one then has
\[
  (v\otimes\lambda)^\star=\lambda^\sharp\otimes v^\flat.
\]
We write the upper odd
layer as \(S^\alpha(\lambda)\) and \(\bar S^{\dot\alpha}(\lambda)\) for
\(\lambda\in E^*\), and we set
\[
  P_{\alpha\dot\beta}:=(\sigma^\mu)_{\alpha\dot\beta}P_\mu,
  \qquad
  K^{\alpha\dot\beta}:=(\bar\sigma^\mu)^{\alpha\dot\beta}K_\mu.
\]
We regard \(Q_\alpha(\cdot)\) and \(S^\alpha(\cdot)\) as linear in their labels.
For the barred families we state all internal-action formulas in the invariant
\(A^\star\)-/\(\flat,\sharp\)-form; the component formulas in an orthonormal basis
are the unambiguous concrete version.

\begin{lemma}[Normalization of the upper odd layer]
\label{lem:4b-upper-odd-normalization}
After choosing bases of \(W_S\), one may and shall normalize the \(S\)-operators
so that on the common core \(\mathcal D_{\mathrm{sc}}\) the unique nonzero
\(\mathfrak{so}(1,3)\)-intertwiners are
\[
  [P_{\alpha\dot\beta},S^\gamma(\lambda)]
  =
  2\,\delta_\alpha^\gamma\,\bar Q_{\dot\beta}(\lambda^\sharp),
  \qquad
  [P_{\alpha\dot\beta},\bar S^{\dot\gamma}(\lambda)]
  =
  2\,\delta_{\dot\beta}^{\dot\gamma}\,Q_\alpha(\lambda^\sharp),
\]
\[
  [K^{\alpha\dot\beta},Q_\gamma(v)]
  =
  2\,\delta_\gamma^\alpha\,\bar S^{\dot\beta}(v^\flat),
  \qquad
  [K^{\alpha\dot\beta},\bar Q_{\dot\gamma}(v)]
  =
  2\,\delta_{\dot\gamma}^{\dot\beta}\,S^\alpha(v^\flat),
\]
with
\[
  [K,S]=[K,\bar S]=0,
  \qquad
  [P,Q]=[P,\bar Q]=0.
\]
\end{lemma}

\begin{proof}
By Theorem~\ref{thm:4b-odd-prolongation}, \(W_S\) is the unique normalized
weight-\(-\tfrac12\) layer in the conformal prolongation of \(W_Q\). For each chirality,
the tensor product of the vector representation with \(W_S\) contains the lower
odd layer with multiplicity one, and likewise \(V\otimes W_Q\) contains \(W_S\)
with multiplicity one. Therefore the corresponding conformal intertwiners are
unique up to an overall nonzero scalar, which we fix once and for all by the
displayed normalization. The annihilation relations by \(P\) and \(K\) are already
part of Theorem~\ref{thm:4b-odd-prolongation}.
\end{proof}

\begin{lemma}[Upper-layer basis may be chosen adjointly]
\label{lem:4b-upper-odd-adjoint-choice}
Assume Theorem~\ref{thm:4b-narrow-common-sector}. The bases of the two chiral
pieces of the upper layer \(W_S\) may be chosen so that
\[
\bigl(S^\alpha(\lambda)\bigr)^\dagger=\bar S^{\dot\alpha}(\lambda)
\]
under the \(\flat/\sharp\) identification, and simultaneously the normalization
conventions of Lemma~\ref{lem:4b-upper-odd-normalization} hold.
\end{lemma}

\begin{proof}
Because \(W_S\) is realized on the same Hilbert space \(\mathcal H\), the Hilbert
adjoint sends one chiral piece of \(W_S\) to the opposite chiral piece. Choose a
basis of one chirality and define the basis of the opposite chirality by taking
adjoints. By Theorem~\ref{thm:odd-conformal-closure-finite}, both chiral pieces
occur with the standard multiplicity-one Lorentz type. Therefore the bridge
intertwiners in Lemma~\ref{lem:4b-upper-odd-normalization} are still unique up to a
common nonzero scalar on each chirality, and that scalar can be fixed once and for
all so that the displayed normalization continues to hold. This preserves the
adjoint pairing.
\end{proof}

\begin{definition}[\texorpdfstring{\(\mathfrak{gl}(2)\oplus\mathfrak{gl}(2)\)}
  {gl(2)+gl(2)} degree-zero generators]
\label{def:4b-degree-zero-generators}
Let
\[
  L_\alpha{}^\beta,
  \qquad
  \bar L^{\dot\alpha}{}_{\dot\beta}
\]
denote the standard complex-linear combinations of \(M_{\mu\nu}\) and \(D\) that
realize the left and right \(\mathfrak{gl}(2)\) actions inside
\(\mathfrak{so}(4,2)_\mathbb C\). They are characterized by
\[
  [L_\alpha{}^\beta,Q_\gamma(v)]
  =
  \delta_\gamma^\beta Q_\alpha(v),
  \qquad
  [L_\alpha{}^\beta,S^\gamma(\lambda)]
  =
  -\delta_\alpha^\gamma S^\beta(\lambda),
\]
\[
  [\bar L^{\dot\alpha}{}_{\dot\beta},\bar Q_{\dot\gamma}(v)]
  =
  \delta_{\dot\gamma}^{\dot\alpha}\bar Q_{\dot\beta}(v),
  \qquad
  [\bar L^{\dot\alpha}{}_{\dot\beta},\bar S^{\dot\gamma}(\lambda)]
  =
  -\delta_{\dot\beta}^{\dot\gamma}\bar S^{\dot\alpha}(\lambda),
\]
and by the analogous tautological action on \(P_{\alpha\dot\beta}\) and
\(K^{\alpha\dot\beta}\). This is only a repackaging of the already available
\(M_{\mu\nu}\) and \(D\).
\end{definition}

\begin{theorem}[Unbarred degree-zero mixed-bracket decomposition]
\label{thm:QS-conformal-plus-internal}
\label{thm:QS-unbarred-first}
\label{lem:QS-conformal-plus-internal}
Assume Theorem~\ref{thm:4b-narrow-common-sector}. Then for every \(v\in E\) and
\(\lambda\in E^*\) there exists a unique even operator \(T_{v,\lambda}\) on
\(\mathcal D_{\mathrm{sc}}\) such that
\[
\{Q_\alpha(v),S^\beta(\lambda)\}
=
\lambda(v)\,L_\alpha{}^\beta+\delta_\alpha^\beta\,T_{v,\lambda},
\]
\[
\{Q_\alpha(v),\bar S^{\dot\beta}(\lambda)\}=0,
\qquad
\{\bar Q_{\dot\alpha}(v),S^\beta(\lambda)\}=0.
\]
Moreover,
\[
  [T_{v,\lambda},P_\mu]
  =
  [T_{v,\lambda},M_{\mu\nu}]
  =
  [T_{v,\lambda},D]
  =0.
\]
\end{theorem}

\begin{proof}
The anticommutator \(\{Q_\alpha(v),S^\beta(\lambda)\}\) is even, has total
\(D\)-weight zero, and transforms in
\[
  S_+\otimes S_+^*
  \cong
  \operatorname{End}(S_+)
  \cong
  \mathfrak{gl}(S_+).
\]
Hence it decomposes uniquely into the tautological \(\mathfrak{gl}(S_+)\) piece
and a scalar-on-spinors remainder; this defines \(T_{v,\lambda}\) and yields the
first displayed formula. The cross-chiral mixed anticommutators have weight zero
and vector Lorentz type, but the degree-zero even conformal algebra contains no
vector piece, so they vanish.

By construction, the scalar remainder commutes with the degree-zero conformal
algebra generated by \(L\) and \(\bar L\). To see that it also commutes with
translations, apply \([P_{\gamma\dot\delta},-]\) to the first mixed-bracket
formula. Lemma~\ref{lem:4b-upper-odd-normalization} and the HLS anticommutator
give
\[
  [P_{\gamma\dot\delta},\{Q_\alpha(v),S^\beta(\lambda)\}]
  =
  4\,\delta_\gamma^\beta\,\lambda(v)\,P_{\alpha\dot\delta}.
\]
With the normalization of Definition~\ref{def:4b-degree-zero-generators}, this is
exactly the translation commutator with
\(\lambda(v)L_\alpha{}^\beta\), so comparison yields
\[
  [P_{\gamma\dot\delta},T_{v,\lambda}]=0.
\]
\end{proof}

\smallskip

\begin{theorem}[Bilinear factorization and faithful lower-layer internal action from mixed brackets]
\label{thm:internal-action-from-mixed-brackets}
Assume Theorem~\ref{thm:QS-conformal-plus-internal} and
Lemma~\ref{lem:HLS-module-injective}. Then:
\begin{enumerate}[(i)]
  \item the assignment \((v,\lambda)\mapsto T_{v,\lambda}\) is bilinear and
        therefore factors uniquely through a linear map
\[
  \Theta:E\otimes E^*\cong\operatorname{End}(E)\longrightarrow
  \operatorname{End}(\mathcal D_{\mathrm{sc}}),
  \qquad
  A\longmapsto T_A,
\]
        with \(T_{v\otimes\lambda}=T_{v,\lambda}\);
  \item for every \(A\in\operatorname{End}(E)\) and every \(w\in E\),
\[
  [T_A,Q_\alpha(w)]
  =
  Q_\alpha(Aw);
\]
  \item the map \(A\mapsto T_A\) is injective.
\end{enumerate}
\end{theorem}

\begin{proof}
The bilinearity of \((v,\lambda)\mapsto T_{v,\lambda}\) is immediate from the
mixed-bracket formula
\[
  \{Q_\alpha(v),S^\beta(\lambda)\}
  =
  \lambda(v)L_\alpha{}^\beta+\delta_\alpha^\beta T_{v,\lambda},
\]
since both the anticommutator and the scalar coefficient \(\lambda(v)\) are
bilinear in \((v,\lambda)\). This gives the unique linear factorization through
\(E\otimes E^*\cong\operatorname{End}(E)\).

For rank-one operators, the same Schur-lemma argument as before shows that
\(T_{v,\lambda}\) acts on the lower odd layer
\[
  W_Q\cong S_+\otimes E.
\]
as the identity on the spinor factor tensored with an endomorphism of \(E\).
Evaluating the super-Jacobi identity on the triple
\((Q_\alpha(v),S^\beta(\lambda),Q_\gamma(w))\) gives
\[
  \lambda(v)\,\delta_\gamma^\beta Q_\alpha(w)
  +
  \delta_\alpha^\beta [T_{v,\lambda},Q_\gamma(w)]
  =
  \lambda(w)\,\delta_\alpha^\beta Q_\gamma(v)
  +
  \delta_\gamma^\beta [T_{w,\lambda},Q_\alpha(v)].
\]
Comparing the \(\delta_\alpha^\beta\) and \(\delta_\gamma^\beta\) components yields
\[
  [T_{v,\lambda},Q_\alpha(w)]
  =
  \lambda(w)\,Q_\alpha(v)
  =
  Q_\alpha\bigl((v\otimes\lambda)w\bigr).
\]
By linearity, this extends to
\[
  [T_A,Q_\alpha(w)]
  =
  Q_\alpha(Aw)
\]
for every \(A\in\operatorname{End}(E)\).

Finally, if \(T_A=0\), then
\[
  0=[T_A,Q_\alpha(w)]=Q_\alpha(Aw)
\]
for every \(w\in E\). Lemma~\ref{lem:HLS-module-injective} therefore gives
\(Aw=0\) for all \(w\), hence \(A=0\).
\end{proof}

Write
\[
  R_{\mathrm{cl},\mathbb C}
  :=
  \Theta(\operatorname{End}(E))
  =
  \{T_A:A\in\operatorname{End}(E)\}.
\]

\begin{lemma}[Faithfulness of the degree-zero internal block on the odd closure]
\label{lem:4b-degree-zero-faithfulness}
Assume Theorem~\ref{thm:HLS-spinor-completeness},
Theorem~\ref{thm:4b-odd-prolongation}, and
Theorem~\ref{thm:4b-narrow-common-sector},
Lemma~\ref{lem:4b-KQ-scalar}, and
Lemma~\ref{lem:4b-upper-D-normal-form}. Let \(U\) be an even operator on
\(\mathcal D_{\mathrm{sc}}\) such that:
\begin{enumerate}[(i)]
  \item \(U\) preserves the odd closure \(W_{\mathrm{sc}}=W_Q\oplus W_S\) under
        commutator;
  \item \(U\) commutes with \(P_\mu\), \(M_{\mu\nu}\), and \(D\);
  \item \(U\) acts trivially on the lower odd layer, i.e.
\[
  [U,Q_\alpha(v)]
  =
  [U,\bar Q_{\dot\alpha}(v)]
  =
  0
  \qquad
  (v\in E).
\]
\end{enumerate}
Then
\[
  [U,S^\alpha(\lambda)]
  =
  [U,\bar S^{\dot\alpha}(\lambda)]
  =
  0
  \qquad
  (\lambda\in E^*).
\]
In particular, inside the generated degree-zero block, an operator is determined
by its action on \(W_Q\).
\end{lemma}

\begin{proof}
Let \(X\) belong to one chiral component of \(W_S\). Since \(U\) commutes with
translations,
\[
  [P_\mu,[U,X]]
  =
  [[P_\mu,U],X]+[U,[P_\mu,X]]
  =
  [U,[P_\mu,X]].
\]
By Theorem~\ref{thm:4b-odd-prolongation}, translations map the upper layer into
the lower layer, and \(U\) annihilates the lower layer by hypothesis. Hence
\([P_\mu,[U,X]]=0\) for every \(\mu\), so \([U,X]\) is a translation-invariant
odd generator.

By Theorem~\ref{thm:HLS-spinor-completeness}, every translation-invariant odd
generator lies in the HLS lower layer \(W_Q\). By Lemma~\ref{lem:4b-KQ-scalar},
one has \(K_Q=\kappa\,\mathbf 1_E\), so Theorem~\ref{thm:4a-Delta-half} and
Lemma~\ref{lem:4b-upper-D-normal-form} give exact scalar \(D\)-weights:
\[
  [D,Q_\alpha(v)]
  =
  \Bigl(\frac{i}{2}+i\kappa\Bigr)Q_\alpha(v),
  \qquad
  [D,S^\alpha(\lambda)]
  =
  \Bigl(-\frac{i}{2}-i\kappa\Bigr)S^\alpha(\lambda),
\]
\[
  [D,\bar Q_{\dot\alpha}(v)]
  =
  \Bigl(\frac{i}{2}-i\kappa\Bigr)\bar Q_{\dot\alpha}(v),
  \qquad
  [D,\bar S^{\dot\alpha}(\lambda)]
  =
  \Bigl(-\frac{i}{2}+i\kappa\Bigr)\bar S^{\dot\alpha}(\lambda).
\]
Because \(U\) commutes with \(D\), the commutator \([U,X]\) has the same
\(D\)-weight as \(X\). But the lower and upper weights are distinct, so a
translation-invariant element of \(W_Q\) cannot have the \(D\)-weight of an
element of \(W_S\). Hence \([U,X]=0\). The barred case is identical.
\end{proof}

\begin{lemma}[The skew part of the lower-layer \(D\)-action is scalar]
\label{lem:4b-KQ-scalar}
Assume Theorem~\ref{thm:4a-Delta-half},
Theorem~\ref{thm:QS-conformal-plus-internal}, and
Theorem~\ref{thm:internal-action-from-mixed-brackets}. Then
\[
  [K_Q,A]=0
  \qquad
  (A\in\operatorname{End}(E)).
\]
Consequently,
\[
  K_Q=\kappa\,\mathbf 1_E
  \qquad
  \text{for a unique }\kappa\in i\mathbb R.
\]
\end{lemma}

\begin{proof}
Fix \(A\in\operatorname{End}(E)\). By linearity,
Theorem~\ref{thm:QS-conformal-plus-internal} gives
\[
  [D,T_A]=0.
\]
Apply this commutator identity to \(Q_\alpha(w)\) for arbitrary \(w\in E\):
\[
  0
  =
  [D,[T_A,Q_\alpha(w)]]-[T_A,[D,Q_\alpha(w)]].
\]
Using \([T_A,Q_\alpha(w)]=Q_\alpha(Aw)\) and the lower-layer normal form from
Theorem~\ref{thm:4a-Delta-half}, one obtains
\[
  0
  =
  Q_\alpha(K_QAw)-Q_\alpha(AK_Qw).
\]
Lemma~\ref{lem:HLS-module-injective} therefore gives
\[
  (K_QA-AK_Q)w=0
  \qquad
  (w\in E),
\]
hence \([K_Q,A]=0\) for every \(A\in\operatorname{End}(E)\). The only
endomorphisms commuting with all of \(\operatorname{End}(E)\) are the scalars, so
\(K_Q=\kappa\,\mathbf 1_E\). The skew-Hermiticity of \(K_Q\) from
Theorem~\ref{thm:4a-Delta-half} forces \(\kappa\in i\mathbb R\).
\end{proof}

\begin{lemma}[Upper-layer dilatation normal form]
\label{lem:4b-upper-D-normal-form}
Assume Theorem~\ref{thm:4b-narrow-common-sector},
Lemma~\ref{lem:4b-upper-odd-normalization},
Theorem~\ref{thm:4a-Delta-half}, and
Lemma~\ref{lem:HLS-module-injective}. Then there exist endomorphisms
\(B_S,C_S\in\operatorname{End}(E^*)\) such that
\[
  [D,S^\alpha(\lambda)]
  =
  i\,S^\alpha(B_S\lambda),
  \qquad
  [D,\bar S^{\dot\alpha}(\lambda)]
  =
  i\,\bar S^{\dot\alpha}(C_S\lambda),
\]
and in fact
\[
  B_S=-\frac12\mathbf 1_{E^*}-K_Q^\vee,
  \qquad
  C_S=-\frac12\mathbf 1_{E^*}+K_Q^\vee.
\]
Equivalently,
\[
  [D,S^\alpha(\lambda)]
  =
  -\frac{i}{2}S^\alpha(\lambda)-iS^\alpha(K_Q^\vee\lambda),
\]
\[
  [D,\bar S^{\dot\alpha}(\lambda)]
  =
  -\frac{i}{2}\bar S^{\dot\alpha}(\lambda)+i\bar S^{\dot\alpha}(K_Q^\vee\lambda).
\]
\end{lemma}

\begin{proof}
Because \(W_S\) is finite-dimensional, \(D\) preserves \(W_S\), and \(D\)
commutes with the Lorentz action, Schur-type reasoning on the spinor factor gives
endomorphisms \(B_S\) and \(C_S\) with the displayed form.

Apply \(D\) to the bridge relation
\[
  [P_{\beta\dot\gamma},S^\alpha(\lambda)]
  =
  2\,\delta_\beta^\alpha\,\bar Q_{\dot\gamma}(\lambda^\sharp).
\]
Using \([D,P_{\beta\dot\gamma}]=iP_{\beta\dot\gamma}\) and the barred lower-layer
formula from Theorem~\ref{thm:4a-Delta-half}, one gets
\[
\begin{aligned}
[D,[P_{\beta\dot\gamma},S^\alpha(\lambda)]]
&=
[[D,P_{\beta\dot\gamma}],S^\alpha(\lambda)]
+[P_{\beta\dot\gamma},[D,S^\alpha(\lambda)]]\\
&=
i[P_{\beta\dot\gamma},S^\alpha(\lambda)]
+i[P_{\beta\dot\gamma},S^\alpha(B_S\lambda)]\\
&=
2\,\delta_\beta^\alpha\,i\Bigl(
\bar Q_{\dot\gamma}(\lambda^\sharp)+
\bar Q_{\dot\gamma}((B_S\lambda)^\sharp)
\Bigr),
\end{aligned}
\]
while the left-hand side is also
\[
  [D,2\,\delta_\beta^\alpha\,\bar Q_{\dot\gamma}(\lambda^\sharp)]
  =
  2\,\delta_\beta^\alpha\Bigl(
  \frac{i}{2}\bar Q_{\dot\gamma}(\lambda^\sharp)
  -i\bar Q_{\dot\gamma}(K_Q\lambda^\sharp)
  \Bigr).
\]
Cancelling the common factor and using Lemma~\ref{lem:HLS-module-injective}
through the \(\flat/\sharp\) identification gives
\[
  (B_S\lambda)^\sharp
  =
  -\frac12\lambda^\sharp-K_Q\lambda^\sharp,
\]
which is exactly \(B_S=-\tfrac12\mathbf 1_{E^*}-K_Q^\vee\).

The barred formula is obtained identically from the bridge
\[
  [P_{\beta\dot\gamma},\bar S^{\dot\alpha}(\lambda)]
  =
  2\,\delta_{\dot\gamma}^{\dot\alpha}Q_\beta(\lambda^\sharp)
\]
and the unbarred lower-layer formula from Theorem~\ref{thm:4a-Delta-half},
giving \(C_S=-\tfrac12\mathbf 1_{E^*}+K_Q^\vee\).
\end{proof}

\begin{corollary}[Internal action on all four odd families]
\label{cor:4b-contragredient-internal-action}
Assume Theorem~\ref{thm:Main-HLS},
Theorem~\ref{thm:HLS-spinor-completeness},
Theorem~\ref{thm:QS-conformal-plus-internal},
Theorem~\ref{thm:internal-action-from-mixed-brackets},
Lemma~\ref{lem:4b-upper-odd-normalization},
Lemma~\ref{lem:4b-KQ-scalar}, and
Lemma~\ref{lem:4b-upper-D-normal-form}. Then for every
\(A\in\operatorname{End}(E)\), \(v\in E\), and \(\mu\in E^*\),
\[
[T_A,Q_\alpha(v)]
=
Q_\alpha(Av),
\qquad
[T_A,\bar Q_{\dot\alpha}(v)]
=
-\bar Q_{\dot\alpha}(A^\star v),
\]
\[
[T_A,S^\alpha(\mu)]
=
-S^\alpha(A^\vee\mu),
\qquad
[T_A,\bar S^{\dot\alpha}(\mu)]
=
\bar S^{\dot\alpha}\bigl((A\mu^\sharp)^\flat\bigr).
\]
In particular, if \(\{e_I\}\) is an orthonormal basis of \(E\), \(\{e^I\}\) its
dual basis, and \(R^I{}_J:=T_{e_I,e^J}\), then
\[
[R^I{}_J,Q_\alpha^K]=\delta_J^K\,Q_\alpha^I,
\qquad
[R^I{}_J,\bar Q_{\dot\alpha K}]
=
-\delta_K^I\,\bar Q_{\dot\alpha J},
\]
\[
[R^I{}_J,S^\alpha{}_K]
=
-\delta_K^I\,S^\alpha{}_J,
\qquad
[R^I{}_J,\bar S^{\dot\alpha K}]
=
\delta_J^K\,\bar S^{\dot\alpha I}.
\]
\end{corollary}

\begin{proof}
The first formula is exactly
Theorem~\ref{thm:internal-action-from-mixed-brackets}.

For the barred lower layer, fix \(A\in\operatorname{End}(E)\) and \(w\in E\).
Because \([T_A,P_\mu]=0\) by linearity from
Theorem~\ref{thm:QS-conformal-plus-internal}, and because
\([P_\mu,\bar Q_{\dot\alpha}(w)]=0\), the commutator
\[
Y_{\dot\alpha}(w):=[T_A,\bar Q_{\dot\alpha}(w)]
\]
is again a translation-invariant barred odd Weyl-spinor symmetry generator.
By Theorem~\ref{thm:HLS-spinor-completeness}, there exists a unique
\(B_Aw\in E\) such that
\[
Y_{\dot\alpha}(w)=\bar Q_{\dot\alpha}(B_Aw).
\]
Now use the HLS anticommutator and \([T_A,P_{\alpha\dot\beta}]=0\):
\[
\begin{aligned}
0
&=
[T_A,\{Q_\alpha(v),\bar Q_{\dot\beta}(w)\}] \\
&=
\{[T_A,Q_\alpha(v)],\bar Q_{\dot\beta}(w)\}
+
\{Q_\alpha(v),[T_A,\bar Q_{\dot\beta}(w)]\} \\
&=
\{Q_\alpha(Av),\bar Q_{\dot\beta}(w)\}
+
\{Q_\alpha(v),\bar Q_{\dot\beta}(B_Aw)\} \\
&=
2\bigl(
\langle Av,w\rangle_Q+\langle v,B_Aw\rangle_Q
\bigr)P_{\alpha\dot\beta}.
\end{aligned}
\]
Choose \((\alpha,\dot\beta)\) with \(P_{\alpha\dot\beta}\neq 0\), using
Lemma~\ref{lem:4a-HLS-nonzero-translation}. Then
\[
\langle Av,w\rangle_Q+\langle v,B_Aw\rangle_Q=0
\qquad
(v,w\in E),
\]
so by definition of \(A^\star\),
\[
B_A=-A^\star.
\]
Thus
\[
[T_A,\bar Q_{\dot\alpha}(w)]
=
-\bar Q_{\dot\alpha}(A^\star w).
\]

For the unbarred upper layer, fix \(\mu\in E^*\) and define
\[
X^\alpha(\mu)
:=
[T_A,S^\alpha(\mu)]+S^\alpha(A^\vee\mu).
\]
Using \([T_A,P_{\beta\dot\gamma}]=0\) and the bridge relation from
Lemma~\ref{lem:4b-upper-odd-normalization},
\[
[P_{\beta\dot\gamma},S^\alpha(\mu)]
=
2\,\delta_\beta^\alpha\,\bar Q_{\dot\gamma}(\mu^\sharp),
\]
we compute
\[
\begin{aligned}
[P_{\beta\dot\gamma},X^\alpha(\mu)]
&=
[T_A,[P_{\beta\dot\gamma},S^\alpha(\mu)]]
+
[P_{\beta\dot\gamma},S^\alpha(A^\vee\mu)] \\
&=
2\,\delta_\beta^\alpha\,[T_A,\bar Q_{\dot\gamma}(\mu^\sharp)]
+
2\,\delta_\beta^\alpha\,\bar Q_{\dot\gamma}((A^\vee\mu)^\sharp).
\end{aligned}
\]
But \((A^\vee\mu)^\sharp=A^\star\mu^\sharp\), so the barred lower-layer formula
just proved gives
\[
[T_A,\bar Q_{\dot\gamma}(\mu^\sharp)]
=
-\bar Q_{\dot\gamma}(A^\star\mu^\sharp)
=
-\bar Q_{\dot\gamma}((A^\vee\mu)^\sharp).
\]
Hence
\[
[P_{\beta\dot\gamma},X^\alpha(\mu)]=0
\qquad
(\beta,\dot\gamma).
\]
So \(X^\alpha(\mu)\) is translation-invariant. By
Theorem~\ref{thm:HLS-spinor-completeness}, there exists a unique \(u\in E\) with
\[
X^\alpha(\mu)=Q_\alpha(u).
\]

Now use Lemma~\ref{lem:4b-KQ-scalar}, writing \(K_Q=\kappa\,\mathbf 1_E\) with
\(\kappa\in i\mathbb R\). Then
Theorem~\ref{thm:4a-Delta-half} and
Lemma~\ref{lem:4b-upper-D-normal-form} give the scalar \(D\)-weights
\[
[D,Q_\alpha(u)]
=
\Bigl(\frac{i}{2}+i\kappa\Bigr)Q_\alpha(u),
\qquad
[D,S^\alpha(\mu)]
=
\Bigl(-\frac{i}{2}-i\kappa\Bigr)S^\alpha(\mu).
\]
Since \([T_A,D]=0\), we obtain
\[
[D,X^\alpha(\mu)]
=
\Bigl(-\frac{i}{2}-i\kappa\Bigr)X^\alpha(\mu).
\]
But \(X^\alpha(\mu)=Q_\alpha(u)\), so the same operator would have to carry both
eigenvalues \(\frac{i}{2}+i\kappa\) and \(-\frac{i}{2}-i\kappa\), which differ by
\(i\). Therefore \(u=0\), hence \(X^\alpha(\mu)=0\). This proves
\[
[T_A,S^\alpha(\mu)]
=
-S^\alpha(A^\vee\mu).
\]

For the barred upper layer, fix \(\mu\in E^*\) and define
\[
Z^{\dot\alpha}(\mu)
:=
[T_A,\bar S^{\dot\alpha}(\mu)]
-\bar S^{\dot\alpha}\bigl((A\mu^\sharp)^\flat\bigr).
\]
Again using \([T_A,P_{\beta\dot\gamma}]=0\) and the second bridge relation from
Lemma~\ref{lem:4b-upper-odd-normalization},
\[
[P_{\beta\dot\gamma},\bar S^{\dot\alpha}(\mu)]
=
2\,\delta_{\dot\gamma}^{\dot\alpha}\,Q_\beta(\mu^\sharp),
\]
we get
\[
\begin{aligned}
[P_{\beta\dot\gamma},Z^{\dot\alpha}(\mu)]
&=
[T_A,[P_{\beta\dot\gamma},\bar S^{\dot\alpha}(\mu)]]
-
[P_{\beta\dot\gamma},
\bar S^{\dot\alpha}((A\mu^\sharp)^\flat)] \\
&=
2\,\delta_{\dot\gamma}^{\dot\alpha}\,[T_A,Q_\beta(\mu^\sharp)]
-
2\,\delta_{\dot\gamma}^{\dot\alpha}\,Q_\beta(((A\mu^\sharp)^\flat)^\sharp) \\
&=
2\,\delta_{\dot\gamma}^{\dot\alpha}\,Q_\beta(A\mu^\sharp)
-
2\,\delta_{\dot\gamma}^{\dot\alpha}\,Q_\beta(A\mu^\sharp) \\
&=
0.
\end{aligned}
\]
So \(Z^{\dot\alpha}(\mu)\) is translation-invariant. By
Theorem~\ref{thm:HLS-spinor-completeness}, there exists \(u\in E\) such that
\[
Z^{\dot\alpha}(\mu)=\bar Q_{\dot\alpha}(u).
\]
Using again the scalar \(D\)-weights,
\[
[D,\bar Q_{\dot\alpha}(u)]
=
\Bigl(\frac{i}{2}-i\kappa\Bigr)\bar Q_{\dot\alpha}(u),
\qquad
[D,\bar S^{\dot\alpha}(\mu)]
=
\Bigl(-\frac{i}{2}+i\kappa\Bigr)\bar S^{\dot\alpha}(\mu),
\]
and \([T_A,D]=0\), we get
\[
[D,Z^{\dot\alpha}(\mu)]
=
\Bigl(-\frac{i}{2}+i\kappa\Bigr)Z^{\dot\alpha}(\mu).
\]
Again the two eigenvalues differ by \(i\), so \(u=0\). Therefore
\[
[T_A,\bar S^{\dot\alpha}(\mu)]
=
\bar S^{\dot\alpha}\bigl((A\mu^\sharp)^\flat\bigr).
\]

Finally, specializing to \(A=e_I\otimes e^J\) gives the displayed basis formulas.
\end{proof}

\begin{lemma}[Internal \(\star\) on the multiplicity block]
\label{lem:4b-internal-adjoint}
The HLS Hermitian form on \(E\) defines an involution
\[
  A\longmapsto A^\star
  \qquad
  \text{on }\operatorname{End}(E)
\]
by
\[
  \langle Av,w\rangle_Q=\langle v,A^\star w\rangle_Q.
\]
For rank-one operators,
\[
  (v\otimes\lambda)^\star=\lambda^\sharp\otimes v^\flat.
\]
The \(\star\)-skew part is
\[
  \mathfrak u(E)
  :=
  \{A\in\operatorname{End}(E):A^\star=-A\}
  \cong
  \mathfrak u(N),
\]
and its commutator subalgebra is
\[
  \mathfrak{su}(E)\cong\mathfrak{su}(N).
\]
\end{lemma}

\begin{proof}
For rank-one operators one computes
\[
\langle (v\otimes\lambda)w,x\rangle_Q
=
\lambda(w)\,\langle v,x\rangle_Q
=
\langle w,(\lambda^\sharp\otimes v^\flat)x\rangle_Q,
\]
which gives \((v\otimes\lambda)^\star=\lambda^\sharp\otimes v^\flat\). The final
two statements are the standard linear algebra of a finite-dimensional Hermitian
space.
\end{proof}

\begin{corollary}[Correct barred mixed-bracket formula]
\label{cor:4b-barred-mixed-formula}
Assume Theorem~\ref{thm:internal-action-from-mixed-brackets},
Lemma~\ref{lem:4b-degree-zero-faithfulness},
Lemma~\ref{lem:4b-upper-odd-adjoint-choice}, and
Lemma~\ref{lem:4b-internal-adjoint}. Then
\[
\{\bar Q_{\dot\alpha}(v),\bar S^{\dot\beta}(\lambda)\}
=
\lambda(v)\,\bar L^{\dot\beta}{}_{\dot\alpha}
-\delta_{\dot\alpha}^{\dot\beta}\,T_{(v\otimes\lambda)^\star}.
\]
\end{corollary}

\begin{proof}
Write the barred same-chirality mixed anticommutator in the form
\[
\{\bar Q_{\dot\alpha}(v),\bar S^{\dot\beta}(\lambda)\}
=
\lambda(v)\,\bar L^{\dot\beta}{}_{\dot\alpha}
+\delta_{\dot\alpha}^{\dot\beta}\,B_{v,\lambda}.
\]
Apply the super-Jacobi identity to the triple
\((Q_\gamma(w),\bar Q_{\dot\alpha}(v),\bar S^{\dot\beta}(\lambda))\). Because
\(\{Q,\bar S\}=0\), the right-hand side reduces to the \(P\)--\(\bar S\) bridge
term; using the HLS anticommutator and the normalized bridge relation of
Lemma~\ref{lem:4b-upper-odd-normalization} gives
\[
  [B_{v,\lambda},Q_\gamma(w)]
  =
  Q_\gamma\bigl((v\otimes\lambda)^\star w\bigr).
\]
By Theorem~\ref{thm:internal-action-from-mixed-brackets}, the internal operator
\(T_{(v\otimes\lambda)^\star}\) has the same action on the full lower odd layer.
Set
\[
  U:=B_{v,\lambda}+T_{(v\otimes\lambda)^\star}.
\]
Then \(U\) acts trivially on the lower layer. Since \(U\) is degree-zero
internal, it commutes with translations, Lorentz generators, and \(D\); by
Lemma~\ref{lem:4b-degree-zero-faithfulness}, it therefore acts trivially on the
entire odd closure. Thus, inside the generated degree-zero internal block,
\[
  B_{v,\lambda}=-\,T_{(v\otimes\lambda)^\star},
\]
which is the displayed formula.
\end{proof}

\begin{theorem}[Closure-generated internal saturation]
\label{thm:closure-generated-internal-saturation}
Assume Theorem~\ref{thm:internal-action-from-mixed-brackets},
Lemma~\ref{lem:4b-degree-zero-faithfulness}, and
Lemma~\ref{lem:4b-internal-adjoint}. Then the internal degree-zero block is a
faithful copy of \(\operatorname{End}(E)\): for all
\(A,B\in\operatorname{End}(E)\),
\[
  R_{\mathrm{cl},\mathbb C}
  =
  \{T_A:A\in\operatorname{End}(E)\},
  \qquad
  [T_A,T_B]=T_{[A,B]}.
\]
Under this identification, the \(\star\)-skew part is
\(\mathfrak u(E)\cong\mathfrak u(N)\), and its commutator subalgebra is
\(\mathfrak{su}(E)\cong\mathfrak{su}(N)\).
\end{theorem}

\begin{proof}
Theorem~\ref{thm:internal-action-from-mixed-brackets} identifies every rank-one
operator \(v\otimes\lambda\) with \(T_{v,\lambda}\), and such rank-one operators
span \(\operatorname{End}(E)\). Hence
\[
  R_{\mathrm{cl},\mathbb C}
  =
  \{T_A:A\in\operatorname{End}(E)\}.
\]
For the commutator law, let \(A,B\in\operatorname{End}(E)\). Then for every
\(w\in E\),
\[
  [[T_A,T_B],Q_\alpha(w)]
  =
  Q_\alpha([A,B]w)
  =
  [T_{[A,B]},Q_\alpha(w)].
\]
The difference \([T_A,T_B]-T_{[A,B]}\) is again degree-zero internal, so
Lemma~\ref{lem:4b-degree-zero-faithfulness} shows that it vanishes on the odd
closure. Thus \([T_A,T_B]=T_{[A,B]}\) in the generated internal block.

The \(\star\)-skew part and its commutator subalgebra are exactly
\(\mathfrak u(E)\) and \(\mathfrak{su}(E)\) by
Lemma~\ref{lem:4b-internal-adjoint}.
\end{proof}

\begin{theorem}[The \(S\)-\(\bar S\) anticommutator]
\label{thm:4b-SbarS-anticommutator}
Assume Theorem~\ref{thm:4b-odd-prolongation},
Theorem~\ref{thm:QS-conformal-plus-internal},
Corollary~\ref{cor:4b-barred-mixed-formula},
Theorem~\ref{thm:internal-action-from-mixed-brackets}, and
Corollary~\ref{cor:4b-dilatation-normalization}. Then on \(\mathcal D_{\mathrm{sc}}\) one
has
\[
\{S^\alpha(\lambda),\bar S^{\dot\beta}(\mu)\}
=
2\,\langle \mu^\sharp,\lambda^\sharp\rangle_Q\,K^{\alpha\dot\beta},
\]
with the coefficient fixed by the normalization conventions of
Lemma~\ref{lem:4b-upper-odd-normalization}. Moreover,
\[
\{S,S\}=\{\bar S,\bar S\}=0.
\]
\end{theorem}

\begin{proof}
By Corollary~\ref{cor:4b-dilatation-normalization}, the upper odd generators have
exact \(D\)-weight \(-\tfrac12\). Therefore
\(\{S^\alpha(\lambda),\bar S^{\dot\beta}(\mu)\}\) is an even operator of
\(D\)-weight \(-1\) and Lorentz type \((\tfrac12,\tfrac12)\). The internal block
has weight \(0\) and no vector piece, so inside the even block there is only one
possible generator of this type, namely \(K^{\alpha\dot\beta}\). Hence there
exists a bilinear scalar \(c(\lambda,\mu)\) such that
\[
\{S^\alpha(\lambda),\bar S^{\dot\beta}(\mu)\}
=
c(\lambda,\mu)\,K^{\alpha\dot\beta}.
\]

Apply \([P_{\gamma\dot\delta},-]\) to both sides. Using the normalized bridge
relations of Lemma~\ref{lem:4b-upper-odd-normalization} gives
\[
\begin{aligned}
[P_{\gamma\dot\delta},\{S^\alpha(\lambda),\bar S^{\dot\beta}(\mu)\}]
&=
2\delta_\gamma^\alpha
\{\bar Q_{\dot\delta}(\lambda^\sharp),\bar S^{\dot\beta}(\mu)\}
+
2\delta_{\dot\delta}^{\dot\beta}
\{S^\alpha(\lambda),Q_\gamma(\mu^\sharp)\}.
\end{aligned}
\]
Now insert the barred mixed formula and the unbarred mixed formula:
\[
\{\bar Q_{\dot\delta}(\lambda^\sharp),\bar S^{\dot\beta}(\mu)\}
=
\mu(\lambda^\sharp)\,\bar L^{\dot\beta}{}_{\dot\delta}
-\delta_{\dot\delta}^{\dot\beta}\,T_{(\lambda^\sharp\otimes\mu)^\star},
\]
\[
\{S^\alpha(\lambda),Q_\gamma(\mu^\sharp)\}
=
\mu(\lambda^\sharp)\,L_\gamma{}^\alpha
+\delta_\gamma^\alpha\,T_{\mu^\sharp\otimes\lambda}.
\]
But
\[
(\lambda^\sharp\otimes\mu)^\star
=
\mu^\sharp\otimes\lambda,
\]
so the internal terms cancel exactly. Therefore
\[
\begin{aligned}
[P_{\gamma\dot\delta},\{S^\alpha(\lambda),\bar S^{\dot\beta}(\mu)\}]
&=
2\,\langle \mu^\sharp,\lambda^\sharp\rangle_Q
\bigl(
\delta_\gamma^\alpha\bar L^{\dot\beta}{}_{\dot\delta}
+
\delta_{\dot\delta}^{\dot\beta}L_\gamma{}^\alpha
\bigr),
\end{aligned}
\]
and comparison with the conformal commutator
\([P_{\gamma\dot\delta},K^{\alpha\dot\beta}]\) fixes the scalar and yields the
displayed formula.

For same-chirality upper anticommutators, \(\{S,S\}\) and \(\{\bar S,\bar S\}\)
have \(D\)-weight \(-1\) but transform in same-chirality spinor-square
representations rather than vector type, while the internal block has Lorentz
type \((0,0)\). Hence both vanish.
\end{proof}

\begin{corollary}[Internal generators commute with \(K\)]
\label{cor:4b-internal-K-commutation}
Under the hypotheses of Theorem~\ref{thm:4b-SbarS-anticommutator},
\[
[T_A,K_\mu]=0
\qquad
(A\in\operatorname{End}(E)).
\]
Hence the internal operators commute with the full even conformal algebra.
\end{corollary}

\begin{proof}
Fix a rank-one operator \(A=v\otimes\lambda\). Apply
\([K^{\gamma\dot\delta},-]\) to the unbarred mixed-bracket formula
\[
\{Q_\alpha(v),S^\beta(\lambda)\}
=
\lambda(v)L_\alpha{}^\beta+\delta_\alpha^\beta T_A.
\]
Since \([K,S]=0\), the left-hand side becomes
\[
\{[K^{\gamma\dot\delta},Q_\alpha(v)],S^\beta(\lambda)\}
=
2\,\delta_\alpha^\gamma\{\bar S^{\dot\delta}(v^\flat),S^\beta(\lambda)\}.
\]
By Theorem~\ref{thm:4b-SbarS-anticommutator}, this is exactly the conformal
commutator with the \(\lambda(v)L_\alpha{}^\beta\) term, leaving no internal
remainder. Therefore
\[
[K^{\gamma\dot\delta},T_A]=0.
\]
Since the rank-one operators span \(\operatorname{End}(E)\), linearity gives
\([K_\mu,T_A]=0\) for every \(A\in\operatorname{End}(E)\).
\end{proof}

\begin{corollary}[Normalized dilatation has the standard \(K\)-commutator]
\label{cor:4b-dstd-K}
Assume Corollary~\ref{cor:4b-dilatation-normalization} and
Corollary~\ref{cor:4b-internal-K-commutation}. Then
\[
[D_{\mathrm{std}},K_\mu]=-iK_\mu.
\]
\end{corollary}

\begin{proof}
By Definition~\ref{def:2T-conf},
\[
[D,K_\mu]=-iK_\mu.
\]
Since \(R_{K_Q}=T_{K_Q}\) is internal, Corollary~\ref{cor:4b-internal-K-commutation}
gives
\[
[R_{K_Q},K_\mu]=0.
\]
Therefore
\[
[D_{\mathrm{std}},K_\mu]
=
[D-iR_{K_Q},K_\mu]
=
[D,K_\mu]-i[R_{K_Q},K_\mu]
=
-iK_\mu.
\qedhere
\]
\end{proof}

\begin{remark}[Basis form of the internal generators]
Let \(\{e_I\}\) be an orthonormal basis of \(E\), let \(\{e^I\}\) be the dual
basis, and define
\[
  Q_\alpha^I:=Q_\alpha(e_I),
  \qquad
  \bar Q_{\dot\alpha I}:=\bar Q_{\dot\alpha}(e_I),
  \qquad
  S^\alpha{}_I:=S^\alpha(e^I),
  \qquad
  \bar S^{\dot\alpha I}:=\bar S^{\dot\alpha}(e^I),
\]
\[
  R^I{}_J:=T_{e_I,e^J}.
\]
Then
\[
  [R^I{}_J,Q_\alpha^K]=\delta_J^K\,Q_\alpha^I,
  \qquad
  [R^I{}_J,\bar Q_{\dot\alpha K}]=-\delta_K^I\,\bar Q_{\dot\alpha J},
\]
\[
  [R^I{}_J,S^\alpha{}_K]=-\delta_K^I\,S^\alpha{}_J,
  \qquad
  [R^I{}_J,\bar S^{\dot\alpha K}]=\delta_J^K\,\bar S^{\dot\alpha I}.
\]
The trace part
\[
  r:=\sum_I R^I{}_I
\]
spans the \(\mathfrak u(1)\) factor of \(\mathfrak u(N)\). By
Lemma~\ref{lem:4b-KQ-scalar},
\[
  T_{K_Q}=\kappa r,
\]
so the shift \(D_{\mathrm{std}}-D=-i\kappa r\) lies in the central trace
direction.
\end{remark}

\subsection{Full superconformal realization on the same UV Hilbert space}
\label{subsec:4b-superconf}

\paragraph{Normalized dilatation convention.}
From this point onward, once Corollary~\ref{cor:4b-dilatation-normalization},
Corollary~\ref{cor:4b-internal-K-commutation}, and
Corollary~\ref{cor:4b-dstd-K} are available, we write \(D\) for the normalized
operator \(D_{\mathrm{std}}\). Because
\[
  D_{\mathrm{std}}-D=-i\kappa r
\]
is a central internal trace generator, this replacement preserves all conformal
commutators and all commutators with the internal block. It also changes the
degree-zero splitting only by a basis change in the abelian span of the
conformal trace direction and the internal trace direction. Consequently the
block-matrix identification of the final superalgebra is unchanged.

\begin{theorem}[Explicit bracket table on the common core]
\label{thm:4d-SConf-brackets}
Assume Theorem~\ref{thm:4b-odd-prolongation},
Theorem~\ref{thm:4b-narrow-common-sector},
Theorem~\ref{thm:QS-conformal-plus-internal},
Theorem~\ref{thm:internal-action-from-mixed-brackets},
Corollary~\ref{cor:4b-contragredient-internal-action},
Corollary~\ref{cor:4b-barred-mixed-formula},
Theorem~\ref{thm:closure-generated-internal-saturation},
Corollary~\ref{cor:4b-dilatation-normalization},
Theorem~\ref{thm:4b-SbarS-anticommutator}, and
Corollary~\ref{cor:4b-internal-K-commutation},
and Corollary~\ref{cor:4b-dstd-K}. Then on the common
invariant core \(\mathcal D_{\mathrm{sc}}\) the generators
\[
  P_\mu,\ M_{\mu\nu},\ D,\ K_\mu,\ Q_\alpha(v),\ \bar Q_{\dot\alpha}(v),\
  S^\alpha(\lambda),\ \bar S^{\dot\alpha}(\lambda),\ T_{v,\lambda}
\]
satisfy the following relations:
\begin{enumerate}[(i)]
  \item the even conformal generators satisfy the standard conformal relations
        from Definition~\ref{def:2T-conf}, while the internal generators satisfy
\[
  [T_{v,\lambda},T_{w,\mu}]
  =
  T_{\lambda(w)v,\mu}-T_{w,\mu(v)\lambda},
\]
        which is the faithful \(\operatorname{End}(E)\) commutator of the
        internal block \(A\mapsto T_A\), and
\[
  [T_{v,\lambda},P_\mu]
  =
  [T_{v,\lambda},M_{\mu\nu}]
  =
  [T_{v,\lambda},D]
  =
  [T_{v,\lambda},K_\mu]
  =
  0;
\]
  \item the even conformal algebra acts on the odd generators by the standard
        two-layer conformal-spinor representation:
\[
  [P,Q]=0,
  \qquad
  [K,S]=0,
  \qquad
  [D,Q]=\frac{i}{2}Q,
  \qquad
  [D,S]=-\frac{i}{2}S,
\]
        together with the barred analogues and the normalized bridging relations
        of Lemma~\ref{lem:4b-upper-odd-normalization}; the internal algebra acts by
        the rank-one formulas from
        Theorem~\ref{thm:internal-action-from-mixed-brackets} and
        Corollary~\ref{cor:4b-contragredient-internal-action};
  \item the nonzero odd anticommutators are exactly
\[
  \{Q_\alpha(v),\bar Q_{\dot\beta}(w)\}
  =
  2\,\langle v,w\rangle_Q\,P_{\alpha\dot\beta},
\]
\[
  \{S^\alpha(\lambda),\bar S^{\dot\beta}(\mu)\}
  =
  2\,\langle \mu^\sharp,\lambda^\sharp\rangle_Q\,K^{\alpha\dot\beta},
\]
\[
  \{Q_\alpha(v),S^\beta(\lambda)\}
  =
  \lambda(v)\,L_\alpha{}^\beta+\delta_\alpha^\beta\,T_{v,\lambda},
\]
\[
  \{\bar Q_{\dot\alpha}(v),\bar S^{\dot\beta}(\lambda)\}
  =
  \lambda(v)\,\bar L^{\dot\beta}{}_{\dot\alpha}
  -\delta_{\dot\alpha}^{\dot\beta}\,T_{(v\otimes\lambda)^\star},
\]
        while
\[
  \{Q,Q\}=\{\bar Q,\bar Q\}=\{S,S\}=\{\bar S,\bar S\}=0,
  \qquad
  \{Q,\bar S\}=\{\bar Q,S\}=0.
\]
\end{enumerate}
\end{theorem}

\begin{proof}
Item~\textnormal{(i)}: the even conformal commutators are those of
Definition~\ref{def:2T-conf}, now with the normalized generator
\(D=D_{\mathrm{std}}\), with the \(D\)-\(K\) relation supplied by
Corollary~\ref{cor:4b-dstd-K}. The internal degree-zero block is the faithful copy
\(A\mapsto T_A\) of \(\operatorname{End}(E)\) constructed above. For
\(A,B\in\operatorname{End}(E)\), Jacobi on the lower layer gives
\[
  [[T_A,T_B],Q_\alpha(w)]
  =
  Q_\alpha([A,B]w).
\]
By Theorem~\ref{thm:closure-generated-internal-saturation}, this is exactly the
commutator law \([T_A,T_B]=T_{[A,B]}\), whose rank-one form is the displayed
formula. All internal generators commute with \(P_\mu\), \(M_{\mu\nu}\), and the
normalized \(D\) by construction, and with \(K_\mu\) by
Corollary~\ref{cor:4b-internal-K-commutation}.

Item~\textnormal{(ii)}: the lower and upper exact half-weight relations are those
of Corollary~\ref{cor:4b-dilatation-normalization}. The bridge relations are
those of Lemma~\ref{lem:4b-upper-odd-normalization}. The internal action
formulas are the faithful \(A\mapsto T_A\) action on the lower layer together
with the rank-one formulas of
Corollary~\ref{cor:4b-contragredient-internal-action}.

Item~\textnormal{(iii)}: the \(Q\)-\(\bar Q\) relation is the HLS seed relation.
The \(S\)-\(\bar S\) relation is Theorem~\ref{thm:4b-SbarS-anticommutator}. The
unbarred mixed bracket is Theorem~\ref{thm:QS-conformal-plus-internal}, and the
barred mixed bracket is Corollary~\ref{cor:4b-barred-mixed-formula}. The
same-chirality upper anticommutators vanish by
Theorem~\ref{thm:4b-SbarS-anticommutator}, while the cross-chiral mixed
anticommutators vanish by Theorem~\ref{thm:QS-conformal-plus-internal}. Thus
every graded bracket closes on \(\mathcal D_{\mathrm{sc}}\) and no extra odd or
even generators appear.
\end{proof}

\begin{theorem}[Same-\texorpdfstring{$\mathcal H$}{H} realization of the full \(4\)D superconformal algebra]
\label{thm:4d-SConf}
\label{thm:same-H-superconformal-realization}
Assume Theorem~\ref{thm:4d-SConf-brackets}. Then on the same UV Hilbert space
\(\mathcal H\) there exist densely defined operators
\[
P_\mu,\ M_{\mu\nu},\ D,\ K_\mu,\ Q_\alpha^I,\ \bar Q_{\dot\alpha I},\
S^\alpha{}_I,\ \bar S^{\dot\alpha I},\ R^I{}_J
\]
with common invariant core \(\mathcal D_{\mathrm{sc}}\) such that:
\begin{enumerate}[(i)]
  \item the even operators realize the UV conformal algebra together with the
        internal algebra \(\mathfrak u(N)\);
  \item the odd operators realize the two weight layers of the standard \(4\)D
        superconformal odd module, with \(Q,\bar Q\) at weight \(+\tfrac12\) and
        \(S,\bar S\) at weight \(-\tfrac12\);
  \item all graded brackets close on \(\mathcal D_{\mathrm{sc}}\);
  \item the resulting operator Lie superalgebra is the full \(4\)D
        superconformal algebra.
\end{enumerate}
More precisely, by the standard four-dimensional superconformal classification
picture \cite{Nahm1978,Sohnius:1985qm}, for \(N\neq 4\) the resulting algebra is
\(\mathfrak{su}(2,2|N)\), while in the \(N=4\) case the standard simple algebra is
obtained after quotienting the central \(\mathfrak u(1)\), giving
\(\mathfrak{psu}(2,2|4)\).
\end{theorem}

\begin{proof}
Let
\[
  V:=S_+\oplus S_-
  \qquad
  (\dim_\mathbb C V=4).
\]
Let the odd multiplicity space be \(E\). The bracket table gives the
complexified conformal block as the \(V\)-block
\(\mathfrak{sl}(V)\cong\mathfrak{sl}(4,\mathbb C)\), the internal block as the
faithful \(E\)-block \(\operatorname{End}(E)\), and the odd generators as the
off-diagonal blocks
\[
  Q,\bar S\in\operatorname{Hom}(E,V),
  \qquad
  S,\bar Q\in\operatorname{Hom}(V,E).
\]
With the mixed-bracket formulas, the odd-odd anticommutators are
exactly the block compositions whose \(\operatorname{End}(V)\) component gives
the conformal generators and whose \(\operatorname{End}(E)\) component gives the
internal operators \(T_A\). Therefore the complexified generated superalgebra is
the standard block-matrix algebra
\[
  \mathfrak{sl}(V|E)\cong\mathfrak{sl}(4|N).
\]
The real form is the standard one: the conformal block is
\(\mathfrak{su}(2,2)\) by the Hilbert-space adjoint on the conformal generators,
the odd generators pair as \(Q^\dagger=\bar Q\) and \(S^\dagger=\bar S\), and
the internal real block is the \(\star\)-skew part
\(\mathfrak u(E)\cong\mathfrak u(N)\) supplied by the HLS Hermitian form on
\(E\). These are exactly the standard reality conditions for
\(\mathfrak{su}(2,2|N)\). For \(N=4\), the scalar superidentity has zero
supertrace and is central, so quotienting it gives \(\mathfrak{psu}(2,2|4)\).
This proves items~\textnormal{(i)}--\textnormal{(iv)}.
\end{proof}

This establishes the same-\(\mathcal H\) superconformal closure.
